% Qatra — Technical and Strategic Dossier
%
% Compile with XeLaTeX (needed for the Arabic in the title and app names):
%   xelatex -output-directory=build docs/qatra-dossier.tex
%
% Everything factual in here is drawn from the repository, from the live
% Supabase projects, or from the sources listed in the appendix. Figures that
% come from a test run or a query are labelled as such, and the date they were
% taken is on the title page — a dossier that cannot be re-checked is a
% brochure.
\documentclass[11pt,a4paper]{article}

\usepackage{fontspec}
\usepackage[margin=2.4cm]{geometry}
\usepackage{xcolor}
\usepackage{booktabs}
\usepackage{tabularx}
\usepackage{enumitem}
\usepackage{titlesec}
\usepackage{fancyhdr}
\usepackage[hidelinks]{hyperref}
\usepackage{parskip}

\definecolor{qatraRed}{HTML}{E5484D}
\definecolor{qatraNavy}{HTML}{0B2432}
\definecolor{qatraTeal}{HTML}{0E8BA8}
\definecolor{qatraGrey}{HTML}{5A6B75}
\definecolor{qatraGreen}{HTML}{12B76A}
\definecolor{qatraAmber}{HTML}{F5871F}
\definecolor{rule}{HTML}{DCE4E7}

\usepackage{polyglossia}
\setmainlanguage{english}
\setotherlanguage{arabic}
\setmainfont{Segoe UI}
\newfontfamily\arabicfont[Script=Arabic]{Segoe UI}
\newcommand{\ar}[1]{\textarabic{#1}}

\titleformat{\section}{\Large\bfseries\color{qatraNavy}}{\thesection.}{0.6em}{}
\titleformat{\subsection}{\large\bfseries\color{qatraTeal}}{\thesubsection}{0.6em}{}
\titlespacing*{\section}{0pt}{1.4em}{0.6em}

\pagestyle{fancy}
\fancyhf{}
\fancyhead[L]{\small\color{qatraGrey}Qatra — Technical and Strategic Dossier}
\fancyhead[R]{\small\color{qatraGrey}\thepage}
\renewcommand{\headrulewidth}{0.4pt}

\newcommand{\note}[1]{{\small\color{qatraGrey}\textit{#1}}}
\newcommand{\good}{{\color{qatraGreen}\textbf{Built}}}
\newcommand{\open}{{\color{qatraAmber}\textbf{Open}}}
\newcommand{\none}{{\color{qatraRed}\textbf{Not built}}}

\begin{document}

\begin{titlepage}
\centering
\vspace*{3cm}
{\Huge\bfseries\color{qatraNavy} Qatra \ar{قطرة}}\\[0.5em]
{\Large\color{qatraRed} Technical and Strategic Dossier}\\[2em]
{\large Blood-donation matching for Algeria}\\[0.4em]
{\color{qatraGrey} Families ask, associations vouch, donors answer.}\\[3em]

\begin{minipage}{0.8\textwidth}
\centering\small\color{qatraGrey}
Every figure in this document was taken from the repository, from the live
database, or from a named public source. Test results are from runs on the
date below. Nothing here is aspirational unless it appears in
Section~\ref{sec:notbuilt}, which exists to say what is missing.
\end{minipage}

\vfill
{\color{qatraGrey}Yacine Omar Fodhil \\ 24 August 2026}
\end{titlepage}

\tableofcontents
\newpage

% ===========================================================================
\section{Executive summary}

Algeria is the best country in Africa at collecting blood and still runs short
of it. The Agence Nationale du Sang reports roughly \textbf{677\,000 bags} from
over \textbf{800\,000 donors}, a rate of \textbf{13.98 donors per 1\,000
inhabitants} — the highest on the continent for a decade — and describes its own
problem not as a shortage of willing people but as
\emph{``ce déséquilibre entre l'offre et la demande''}: a timing failure. Only
\textbf{22\% of donors give regularly}, which the agency calls
\emph{``relativement faible''}, and its stated request is that citizens
\textbf{register on regular-donor lists}.

Qatra answers that request. It knows the date every donor becomes eligible
again, and it reaches the compatible ones the moment somebody nearby needs
their type.

Three things distinguish it from what already exists:

\begin{enumerate}[leftmargin=1.4em]
  \item \textbf{Requests are vouched for.} A verified association in the same
        wilaya can put its name to a request. The market-leading competitor
        states in its own store listing that it verifies nothing.
  \item \textbf{Health data is handled lawfully.} Consent is recorded per
        purpose with the version of the text shown; a donor's phone number is
        masked until deliberately opened, and every opening is logged in a
        record the donor can read and nobody can delete.
  \item \textbf{The loop closes.} Post, notify, respond, and the family sees
        who is coming. A registry answers \emph{who could give}. Qatra answers
        \emph{who is coming, right now}.
\end{enumerate}

% ===========================================================================
\section{The problem, with evidence}

\subsection{Algeria does not have a donor shortage}

\begin{center}
\begin{tabularx}{\textwidth}{lX}
\toprule
\textbf{Figure} & \textbf{Meaning} \\
\midrule
13.98 per 1\,000 & Donation rate. First in Africa for around a decade. \\
677\,000 bags & Collected, from 800\,000+ donors; up 6.75\% since 2023. \\
\textbf{22\%} & \textbf{Share of donors who give regularly.} ANS: \emph{``relativement faible''}. \\
58 & Wilayas. Blood is coordinated locally; a donor two provinces away cannot help today. \\
\bottomrule
\end{tabularx}
\end{center}

The ANS names the summer as its pressure point — road-accident season — and has
run joint campaigns with the national police and Algérie Télécom to cover it.
Its ask of the public is continuity, not volume:
\emph{``participer à l'alimentation continue de la banque du sang''}.

\subsection{Which means the problem is reach, not recruitment}

Seventy-eight per cent of Algerian donors give once and are never effectively
called again. Nothing in the existing landscape knows when a given donor became
eligible again, or tells them that somebody two kilometres away needs exactly
their type this evening. That is the gap this product occupies.

% ===========================================================================
\section{Competitive landscape}

\note{Install counts and ratings retrieved from Google Play on 24 August 2026.
Platform claims verified against the operators' own sites on the same date.}

\begin{center}
\begin{tabularx}{\textwidth}{@{}lllX@{}}
\toprule
\textbf{Product} & \textbf{Installs} & \textbf{Rating} & \textbf{Observation} \\
\midrule
\ar{بنك الدم الجزائري} (Nabil\_Soft) & 10K+ & 4.7 / 130 & Market leader. \textbf{Contains ads.} \\
\ar{بنك الدم الجزائري} (SangDz) & 1K+ & — & Solo developer based in France. \\
Sadaka \ar{صدقة} & 1K+ & — & — \\
Nabdz & — & — & Best-designed. 2 requests, 3 donors live. \\
ONSC / Ministry platforms & — & — & Announced June 2025. No public address. \\
\bottomrule
\end{tabularx}
\end{center}

Approximately \textbf{12\,000 installs across every blood application in
Algeria}, in a country of 45 million with the continent's highest donation rate.
The field is not crowded; it is empty.

\subsection{The announced national platforms have no front door}

The National Observatory of Civil Society launched a donor platform with the
Algerian Red Crescent on 14 June 2025; the Ministry of Health announced its own
three days later. Fourteen months on:

\begin{itemize}[leftmargin=1.4em]
  \item \texttt{marsad.dz} is the Observatory's institutional site — governance,
        compliance, tenders. It is actively maintained, with content through
        April 2026, and carries no donor platform.
  \item Neither announcement names a URL, an application, or any route by which
        a citizen reaches the system.
  \item The Ministry's own site lists five other digital services it does ship
        — a pharmacy system, a health map, a professional space, a training
        registry — and no blood platform.
\end{itemize}

A launch with ministers, a telecom sponsor and no discoverable product is a
press event. That may change, and the strategic response is not to compete with
it (Section~\ref{sec:cra}).

\subsection{Three findings about the market leader}

\textbf{It carries advertising.} \ar{بنك الدم الجزائري} is listed as
\emph{Contains ads}. Its developer's other titles are \emph{Easy Hairstyles for
Girls}, \emph{PermisDZ}, and several baccalaureate revision apps. Blood
donation is the sixth item in a portfolio.

\textbf{Its data-safety declaration states ``No data collected.''} The
application collects, by its own description, a donor's name, wilaya,
\textbf{blood type} and telephone number. Blood type is health data. That
declaration is inaccurate on its face.

\textbf{It disclaims verification in writing:}

\begin{quote}\itshape
``This application does not provide any medical services or health advice and
does not represent any medical or governmental entity\ldots{} The information
provided is provided by the users.''
\end{quote}

The fraud problem is therefore not an unrecognised gap in the market. It is a
term of service.

\subsection{And the most serious competitor demonstrates the same emptiness}

Nabdz is the best-built of them — 58 wilayas, urgency tiers, a pledge mechanic
comparable to Qatra's response counter. Its live public grid currently shows two
requests and three donors, one of the requests titled \textbf{``Test Patient''}.

% ===========================================================================
\section{Regulatory position}

\note{This section states the reasoning behind decisions already implemented in
the codebase. It is engineering rationale, not legal advice, and the open
question in Section~\ref{sec:residency} should be settled with counsel or with
the ANS before any public launch.}

\subsection{Health data — Loi 18-07 and Loi 25-11}

A patient's name attached to a blood type, and a donor's telephone number
attached to theirs, are health data. Qatra treats them accordingly:

\begin{itemize}[leftmargin=1.4em]
  \item \textbf{Consent is recorded per purpose, with the version of the text
        shown.} Two purposes exist — \texttt{health\_data} and
        \texttt{contact\_sharing} — each stamped with the wording the user
        actually read (\texttt{health-data-v1}, \texttt{contact-sharing-v1}).
        Consent given against superseded wording stops granting access until
        the user agrees again.
  \item \textbf{Consent records have no delete policy.} They are the evidence
        that processing was lawful, and evidence a controller can erase is not
        evidence.
  \item \textbf{Donor numbers are masked until deliberately opened.} Search
        returns \texttt{05 •• •• •• 56}. A single function returns the whole
        number, and it writes a row while doing so.
  \item \textbf{Every reveal is logged, and the donor can read the log.} The
        record names who opened the number, for which donor, on behalf of which
        association, and when. It cannot be updated or deleted by anyone.
  \item \textbf{Data-subject rights have a route in the product} — export,
        correction, deletion — rather than an email address on a policy page.
\end{itemize}

\subsection{Who may vouch — Loi 12-06 on associations}

Verification publishes an attestation under an association's name. An
association is a legal person acting through its statutory representatives, not
through whichever member happens to be holding a phone; if a vouched request
proves fraudulent, the association carries it. Verification is therefore
restricted to association \emph{administrators} in the request's own wilaya,
enforced in the database rather than the interface.

Volunteers keep everything else — the request list, donor search, the work of
mobilising people. They simply cannot sign the association's name.

\note{``Administrator'' is a proxy for ``entitled to bind the association''. A
real committee's mandate structure may be richer, and blood-request vetting is
not itself a regulated medical act. This mapping is marked in the migration and
should be confirmed with the ANS.}

\subsection{No payment, ever}

Blood donation in Algeria is voluntary and unpaid. Paying donors — in cash,
vouchers, credit, priority, or anything exchangeable for them — compromises the
safety of the supply, because a donor with a financial motive has a reason to
conceal a disqualifying history.

This is a documented invariant at the top of the compliance module and a
code-review gate, not a preference. Non-transferable recognition (streaks,
certificates) is permitted precisely because it cannot be exchanged.

\subsection{Open question: data residency}\label{sec:residency}

The production database is hosted in AWS \texttt{us-east-1}. Algerian patients'
names, blood types and telephone numbers therefore rest on servers outside
Algeria. Four files in the repository carry a \texttt{TODO(compliance)} marker
recording this. It is a launch gate, and it is the question a well-informed
panel asks first.

The answer is that nothing prevents it being solved, and this can be checked
rather than taken on trust. The backend is standard PostgreSQL. Across all
twenty-six migrations there is not a single \texttt{create extension}
statement, and no dependency on any proprietary service --- no hosted-only
networking, secrets or storage calls. The one platform-specific thing the
schema uses is \texttt{auth.uid()}, which belongs to GoTrue, an open-source
component that runs the same way in a self-hosted deployment. The client is
handed its address at startup by injection rather than reading it from a
bundler's environment, so pointing the application at an instance hosted in
Algeria is a configuration change and a migration run, not a rewrite.

Two pieces would still need attention on that path, and are named here rather
than discovered later: the push-notification worker is a Deno function and
would need somewhere to run, and the demonstration OTP provider must be
replaced by a real SMS route before any number is trusted.

The distinction worth drawing is between a system that has not yet moved and
one that cannot. A product built on a proprietary mobile backend cannot be
hosted domestically at all, whatever its owner decides. This one is a Postgres
database and a static web application; where it runs is an operational choice
that can be made by whoever ends up accountable for the data.

% ===========================================================================
\section{How the loop closes}

\begin{enumerate}[leftmargin=1.4em]
  \item \textbf{A family posts a request.} Three steps: who needs blood, where,
        how urgent. The hospital is free text — a family knows their hospital by
        name, not by an identifier in a directory.
  \item \textbf{Phone verification interrupts \emph{after} the form, not
        before.} A plea written at 3am is captured while the person still has
        the words; the draft is held and posted the moment the number is
        confirmed. It is never asked for twice.
  \item \textbf{Compatible donors nearby are notified.} Same wilaya, compatible
        blood type, eligible now — not inside the 90-day cooldown — and never
        the person who posted it.
  \item \textbf{A committee may vouch for it.} The console groups requests by
        what to do next: needs review, open a long time, already vouched.
  \item \textbf{A donor responds, and the family sees it.} The count is public
        so twenty people do not arrive for two units; the identities are not.
  \item \textbf{A donor who cannot go says so.} Withdrawal is a state change the
        family sees, not an erasure.
\end{enumerate}

% ===========================================================================
\section{Technical architecture}

\subsection{Stack}

\begin{center}
\begin{tabularx}{\textwidth}{@{}lX@{}}
\toprule
Web application & React 18, Vite 6, TypeScript, Tailwind for layout \\
Mobile shells & Capacitor (Android, iOS) over the same source \\
Backend & Supabase: PostgreSQL, PostgREST, GoTrue, Edge Functions (Deno) \\
Notifications & Web Push (VAPID), queued in the database, sent by a worker \\
Maps & Leaflet over OpenStreetMap \\
Languages & English, French, Arabic, with full right-to-left support \\
\bottomrule
\end{tabularx}
\end{center}

The core package is framework-agnostic by design: configuration is injected
rather than read from the bundler's environment, so a non-Vite host can consume
the same logic.

\subsection{Security model}

Authorisation lives in the database, not the interface. \textbf{54 row-level
security policies} and \textbf{23 \texttt{SECURITY DEFINER} functions} decide
who may read and write what; the client can only ask.

Three principles are worth stating because they were arrived at by correcting
mistakes:

\begin{itemize}[leftmargin=1.4em]
  \item \textbf{One right per predicate.} A single function once gated
        vouching, patient-record access \emph{and} donor search. Narrowing it
        for one purpose silently removed a right unrelated to the decision.
        They are now separate.
  \item \textbf{Two layers, not one.} Revoking a function from
        \texttt{public} does not remove Supabase's separate grants to the
        \texttt{anon} and \texttt{authenticated} roles. Several functions were
        therefore reachable by anonymous callers and stopped only by their own
        internal guard. The test suite now asserts the grants themselves.
  \item \textbf{Logs that cannot be rewritten.} Consent records and contact
        reveals have no update or delete policy at all.
\end{itemize}

\subsection{Notifications}

Database triggers write to an outbox; an edge function drains it. An HTTP call
inside the transaction would make \emph{posting a plea for blood} fail whenever
a push service was slow, and a failed send would leave nothing to retry.

Each claim takes a five-minute lease. \texttt{SKIP LOCKED} alone separates only
simultaneous workers — a second worker a moment later would send the same
notification again, and the test suite caught exactly that.

The payload carries blood type and wilaya and nothing else. A push lands on a
lock screen that anyone nearby can read, so it says only what the public request
list already shows.

\subsection{Blood compatibility}

The full eight-by-eight red-cell table is written out rather than derived, so a
reader can check it against a transfusion chart line by line. Unknown or missing
types answer \emph{false}, never \emph{true}: the safe answer to ``do we know
this is safe?'' is no.

The unit tests check the table from both directions against an independently
written chart, verify that rhesus is not symmetric, and confirm that a donor
whose type is unrecorded is asked rather than refused.

\subsection{Verification and testing}

\begin{center}
\begin{tabularx}{\textwidth}{@{}lrX@{}}
\toprule
\textbf{Suite} & \textbf{Count} & \textbf{What it covers} \\
\midrule
\texttt{verify:db} & 188 & RLS against real PostgreSQL, as several users \\
\texttt{test:unit} & 14 & Blood compatibility both directions; pledge arithmetic \\
Playwright & 48 & End-to-end journeys, three languages, two viewports \\
\midrule
Migrations & 26 & Applied in filename order \\
Screens & 24 & \\
Components & 41 & \\
Translated strings & 388 & Each in all three languages, enforced by the type system \\
\bottomrule
\end{tabularx}
\end{center}

A missing translation is a compilation error, not a visible gap: the string
table is a TypeScript interface, so a key added in one language and forgotten in
another fails the build.

% ===========================================================================
\section{What is built, and what is not}\label{sec:notbuilt}

\begin{center}
\begin{tabularx}{\textwidth}{@{}lX@{}}
\toprule
\textbf{State} & \textbf{Capability} \\
\midrule
\good & Patient/family posts a request; association verifies it \\
\good & Donor discovery, map, urgency and recency ordering \\
\good & Donor responds; family sees the count; withdrawal supported \\
\good & Web push: compatible donors in-wilaya, and the family on a response \\
\good & Consent, contact-reveal logging, data-subject requests \\
\good & Blood compatibility with tests \\
\good & Trilingual EN / FR / AR with right-to-left \\
\good & Committee invite links: a committee brings the donors it already has \\
\good & Public demonstration deployment \\
\midrule
\open & Demonstration OTP is enabled: any number can be claimed \\
\open & Data residency (Section~\ref{sec:residency}) \\
\open & ``Administrator'' as proxy for authority to bind an association \\
\midrule
\none & Native push for the Capacitor Android build (needs FCM) \\
\none & Android build refreshed against the current application \\
\none & Pagination beyond a 200-request cap per wilaya \\
\none & Any integration with a national registry or hospital system \\
\bottomrule
\end{tabularx}
\end{center}

% ===========================================================================
\section{Positioning}\label{sec:cra}

\subsection{Not ``first''}

There is a national registry announced and at least three applications live.
Claiming to be first is both false and checkable in thirty seconds, and
volunteering the competition first reads as command where being caught by it
reads as naivety.

\subsection{The line}

\begin{quote}\itshape
Algeria has a national donor registry. It answers \emph{who could give}. Qatra
answers \emph{who is coming, right now, and can this request be trusted} — and
it does so to a standard a committee can put its name on.
\end{quote}

\subsection{Why the Red Crescent partnership strengthens}

The Algerian Red Crescent is already the state's chosen partner in this space,
having co-launched the June 2025 platform. A CRA wilaya committee vouching
through Qatra is therefore not a rival to the national effort but the local
layer it does not have. The compliance architecture is what makes that
partnership possible: an institution can only lend its name to a system that
handles health data properly.

% ===========================================================================

% ===========================================================================
\section{Presenting this}\label{sec:pitch}

\note{This section is written for the person presenting, not for the reader
being presented to. Everything in it is drawn from the sections above; nothing
new is claimed here.}

\subsection{The sentence to open and close with}

\begin{quote}
\textbf{Algeria is the best country in Africa at collecting blood and still
runs short of it. That is a timing problem, not a willingness problem — and
timing is a software problem.}
\end{quote}

Say it first. Say it again at the end. Everything between the two is evidence
for it.

\subsection{The first thirty seconds}

Open with the fact that contradicts the room's assumption. Everyone expects a
plea about apathy; the ANS's own numbers say the opposite.

\begin{enumerate}[leftmargin=1.4em]
  \item \textbf{13.98 donors per 1\,000 inhabitants.} First in Africa, for
        about a decade. Algerians give.
  \item \textbf{Only 22\% give regularly.} The agency's own word for this is
        \emph{``relativement faible''}.
  \item \textbf{So 78\% gave once and were never effectively called again.}
        Nothing in the country knows the date a given donor became eligible
        again.
\end{enumerate}

Then the turn: \emph{the ANS's stated request to citizens is to register on
regular-donor lists. This is that register, and it makes the call.}

\subsection{The demonstration, in this order}

Four minutes. Resist showing the whole app; show the loop closing, because that
is the thing no competitor has.

\begin{center}
\begin{tabularx}{\textwidth}{@{}rlX@{}}
\toprule
& \textbf{Show} & \textbf{Say} \\
\midrule
1 & A family posts a request & ``Three fields. The hospital is free text,
because a family knows its hospital by name, not by an identifier.'' \\
2 & A committee vouches for it & ``An association in the same wilaya puts its
name to this. The market leader states in its own store listing that it
verifies nothing.'' \\
3 & The badge on the donor's screen & ``And it survives a WhatsApp forward —
the committee is named in the shared message, so a forwarded plea stays
checkable.'' \\
4 & A donor responds & ``Compatible donors in that wilaya were notified. Not
everyone: the right blood types, eligible today, never the person who posted
it.'' \\
5 & The progress bar & ``Three of four units pledged. The twentieth donor can
see the request is covered and go find another one.'' \\
6 & The invite screen & ``And this is how a Red Crescent committee brings the
donor list it already has — a link, never an upload. Each donor signs up and
consents themselves.'' \\
\bottomrule
\end{tabularx}
\end{center}

\subsection{The argument that matters to a public body}

A ministry does not adopt the best application. It adopts the one it
\emph{can} adopt. Frame the competition as a question of eligibility rather
than of features:

\begin{itemize}[leftmargin=1.4em]
  \item The market leader has ten thousand installs and \textbf{publishes
        donors' telephone numbers}, while declaring ``No data collected'' on a
        listing that collects blood type and phone number. No Algerian public
        body can put its name to that.
  \item A product built on a proprietary foreign mobile backend \textbf{cannot
        be hosted in Algeria at all}, whatever its owner decides.
  \item The announced national platforms have no public front door.
\end{itemize}

So the claim is not ``ours is nicer''. It is:

\begin{quote}
\textbf{This is the only one of these that could lawfully become national
infrastructure.}
\end{quote}

\subsection{Three things to run, rather than describe}

Claims are discounted; commands are not. Each of these takes under two minutes.

\begin{center}
\begin{tabularx}{\textwidth}{@{}lX@{}}
\toprule
\textbf{Command} & \textbf{What it proves in front of them} \\
\midrule
\texttt{npm run verify:db} & 188 checks against a real PostgreSQL that is
\emph{not} Supabase. This is the answer to ``can this be hosted in Algeria'' —
the schema runs on any Postgres, on any server, in any country. \\
\texttt{npm run test:unit} & The blood-compatibility table, checked in both
directions against an independently written chart. The one place where being
wrong is a medical error. \\
\texttt{npm test} & 48 end-to-end journeys through a real browser, in three
languages and two viewports. \\
\bottomrule
\end{tabularx}
\end{center}

\subsection{Questions that will be asked, and honest answers}

The weak answers are the ones given defensively. Each of these is better
answered plainly.

\begin{description}[leftmargin=0em, style=unboxed]

\item[``How many people use it?''] \emph{Nobody yet.} Say it without
flinching, then say what it is designed for: one wilaya committee with the
donors it already has. The invite mechanism exists for exactly that, and it
works. A pilot with one committee is the ask, not a claim of traction.

\item[``Where is the data?''] Today, AWS \texttt{us-east-1} — outside Algeria,
and a launch gate rather than a demonstration gate. Then Section
\ref{sec:residency}: no proprietary extensions, no hosted-only services, the
address injected at startup. Offer to run \texttt{verify:db} on the spot.

\item[``How do you know a request is real?''] Three layers, and none of them
is a claim of certainty: the poster's phone number is verified, an association
in the same wilaya can vouch under its own name, and the committee sees a
plausibility panel before it does. Qatra reduces the cost of a lie; it does not
pretend to eliminate it.

\item[``What stops someone harvesting donors' numbers?''] Search returns masked
numbers. Opening one is a deliberate act, restricted to members of an
association in that wilaya, and every opening is written to a log the donor can
read and nobody can delete.

\item[``Isn't this just a WhatsApp group?''] A chain message is
unattributable and never closes. This one carries the vouching committee's name
into the forward, and the family sees who is actually coming.

\item[``What if the state builds its own?''] It announced platforms and did not
ship a front door. And the architecture is deliberately portable — if the right
answer is that the state operates this, that is a handover, not a competitor.

\item[``How does it make money?''] Be careful and be honest. Donor payment is
excluded by law and by design, so there is no consumer revenue here and there
should not be. The realistic routes are institutional: tooling paid for by the
bodies that coordinate donation, or public funding. \textbf{This is not
solved, and claiming otherwise in front of people who fund things is the
fastest way to lose them.}

\item[``Is phone verification real?''] In the demonstration build, no — a
fixed code is accepted, because there is no SMS provider connected. It is a
flag, it is off in production, and it is listed in Section
\ref{sec:notbuilt}. Volunteering this before being caught by it is worth more
than the point it costs.

\end{description}

\subsection{What not to claim}

\begin{itemize}[leftmargin=1.4em]
  \item \textbf{Do not claim users, pilots or partnerships that do not exist.}
        The Red Crescent is the intended partner, not a signed one. One
        checkable exaggeration discredits every figure in this document.
  \item \textbf{Do not claim lives saved.} The application deliberately removed
        that number because nothing counts it.
  \item \textbf{Do not promise medical outcomes.} Qatra moves information
        between people; it does not transfuse anyone.
  \item \textbf{Do not hide the open items.} Section \ref{sec:notbuilt} exists
        so that a panel finds nothing you did not already say. A gap you
        disclosed is competence; the same gap discovered is a credibility
        problem.
\end{itemize}

\subsection{The ask}

A presentation without a request is a report. Ask for the two things that are
actually needed and cannot be built:

\begin{enumerate}[leftmargin=1.4em]
  \item \textbf{One wilaya committee, for one pilot.} Twenty real donors on a
        real list, for one season. That converts every architectural claim in
        this document into evidence.
  \item \textbf{An introduction to the Agence Nationale du Sang.} Two questions
        need an authoritative answer: who may bind an association, and where
        the data must live.
\end{enumerate}

\section{Sources}

\begin{itemize}[leftmargin=1.4em]\small
  \item Agence Nationale du Sang, on regular-donor lists and the summer
        supply–demand imbalance: \url{https://news.radioalgerie.dz/fr/node/47717}
  \item ANS collection figures and donation rate:
        \url{https://www.algerie360.com/don-de-sang-lagence-nationale-donne-lalerte/}
  \item ONSC platform launch:
        \url{https://www.dzair-tube.dz/en/algeria-launches-first-digital-blood-donor-platform-to-strengthen-national-lifeline/}
  \item Ministry of Health platform:
        \url{https://meatechwatch.com/2025/06/18/algeria-launches-digital-platform-to-enhance-blood-donation-management-nationwide/}
  \item National Observatory of Civil Society: \url{https://marsad.dz}
  \item Ministry of Health: \url{https://www.sante.gov.dz}
  \item \ar{بنك الدم الجزائري} (Nabil\_Soft) on Google Play:
        \url{https://play.google.com/store/apps/details?id=com.dz.bank.blood.nabilsoft.bankblooddz}
  \item \ar{بنك الدم الجزائري} SangDz on Google Play:
        \url{https://play.google.com/store/apps/details?id=com.blood.sangdz}
  \item Nabdz emergency grid: \url{https://nabdz.com/en/blood-donation}
\end{itemize}

\vfill
\note{Repository: \url{https://github.com/2him29/startup-idea} — public
demonstration at \url{https://2him29.github.io/startup-idea/}}

\end{document}
